\documentclass{beamer}
\usepackage[latin1]{inputenc}
\beamertemplateshadingbackground{red!10}{blue!10}
%\usepackage{fancybox}
\usepackage{epsfig}
\usepackage{verbatim}
\usepackage{url}
%\usepackage{graphics}
%\usepackage{xcolor}
\usepackage{fancybox}
\usepackage{moreverb}
%\usepackage[all]{xy}
\usepackage{listings}
\usepackage{filecontents}
\usepackage{graphicx}

\lstset{
  language=Lisp,
  basicstyle=\scriptsize\ttfamily,
  keywordstyle={},
  commentstyle={},
  stringstyle={}}

\def\inputfig#1{\input #1}
\def\inputeps#1{\includegraphics{#1}}
\def\inputtex#1{\input #1}

\inputtex{logos.tex}

%\definecolor{ORANGE}{named}{Orange}

\definecolor{GREEN}{rgb}{0,0.8,0}
\definecolor{YELLOW}{rgb}{1,1,0}
\definecolor{ORANGE}{rgb}{1,0.647,0}
\definecolor{PURPLE}{rgb}{0.627,0.126,0.941}
\definecolor{PURPLE}{named}{purple}
\definecolor{PINK}{rgb}{1,0.412,0.706}
\definecolor{WHEAT}{rgb}{1,0.8,0.6}
\definecolor{BLUE}{rgb}{0,0,1}
\definecolor{GRAY}{named}{gray}
\definecolor{CYAN}{named}{cyan}

\newcommand{\orchid}[1]{\textcolor{Orchid}{#1}}
\newcommand{\defun}[1]{\orchid{#1}}

\newcommand{\BROWN}[1]{\textcolor{BROWN}{#1}}
\newcommand{\RED}[1]{\textcolor{red}{#1}}
\newcommand{\YELLOW}[1]{\textcolor{YELLOW}{#1}}
\newcommand{\PINK}[1]{\textcolor{PINK}{#1}}
\newcommand{\WHEAT}[1]{\textcolor{wheat}{#1}}
\newcommand{\GREEN}[1]{\textcolor{GREEN}{#1}}
\newcommand{\PURPLE}[1]{\textcolor{PURPLE}{#1}}
\newcommand{\BLACK}[1]{\textcolor{black}{#1}}
\newcommand{\WHITE}[1]{\textcolor{WHITE}{#1}}
\newcommand{\MAGENTA}[1]{\textcolor{MAGENTA}{#1}}
\newcommand{\ORANGE}[1]{\textcolor{ORANGE}{#1}}
\newcommand{\BLUE}[1]{\textcolor{BLUE}{#1}}
\newcommand{\GRAY}[1]{\textcolor{gray}{#1}}
\newcommand{\CYAN}[1]{\textcolor{cyan }{#1}}

\newcommand{\reference}[2]{\textcolor{PINK}{[#1~#2]}}
%\newcommand{\vect}[1]{\stackrel{\rightarrow}{#1}}

% Use some nice templates
\beamertemplatetransparentcovereddynamic

\newcommand{\A}{{\mathbb A}}
\newcommand{\degr}{\mathrm{deg}}

\title{Choosing a programming language}

\author{Robert Strandh}
\institute{
University of Bordeaux
}
\date{May, 2018}

%\inputtex{macros.tex}


\begin{document}
\frame{
\resizebox{3cm}{!}{\includegraphics{Logobx.pdf}}
\hfill
\resizebox{1.5cm}{!}{\includegraphics{labri-logo.pdf}}
\titlepage
\vfill
\small{Dfind, G�teborg}
}

\setbeamertemplate{footline}{
\vspace{-1em}
\hspace*{1ex}{~} \GRAY{\insertframenumber/\inserttotalframenumber}
}

\begin{frame}
\frametitle{Overview of talk}
\begin{itemize}
\item Programming language characteristics.
\item Common misconceptions.
\item Requirements for making a good choice.
\end{itemize}
\end{frame}

%-----------------------------------------------------------
\begin{frame}\frametitle{How the choice is often made}

The choice of programming language for a project (when there are
several possibilities) is often based on \emph{gut feeling}.
\vskip 0.5cm
Often, no real analysis is made, because the decision maker:

\begin{itemize}
\item has only partial knowledge of the characteristics of possible
  choices;
\item sometimes has incorrect information about the possible choices;
\item has insufficient training to appreciate the characteristics of
  possible choices;
\end{itemize}

\end{frame}
%-----------------------------------------------------------
\begin{frame}\frametitle{How the choice is often made}

Often, no real analysis is made, because the decision maker:

\begin{itemize}
\item has insufficient experience with the possible choices to
  determine which choice is adapted to the current project;
\item has insufficient information about the cost associated with
  training staff in a new language vs the productivity advantages of
  that language;
\item has insufficient information about the cost associated with
  hiring new staff for a new language vs the productivity advantages of
  that language.
\end{itemize}

\end{frame}
%-----------------------------------------------------------
\begin{frame}\frametitle{Language vs implementation}

\end{frame}
%-----------------------------------------------------------
\begin{frame}\frametitle{Language vs implementation}

Language: A description of the syntax and semantics of conforming
programs, and of consequences of using non-conforming constructs.
\vskip 0.5cm
Example of the latter: In C, obtaining a pointer outside of an array
has undefined consequences.  (It is interesting to contemplate why.)

\end{frame}
%-----------------------------------------------------------
\begin{frame}\frametitle{Language vs implementation}

Implementation: Software that accepts conforming programs and executes
them according to the language semantics, and that reports
non-conforming constructs where required by the language definition.

\end{frame}
%-----------------------------------------------------------
\begin{frame}\frametitle{Language vs implementation}

In many cases, the language definition does not require the compiler
to check for non-conforming constructs.  Why?
\vskip 0.5cm
As a consequence, many non-conforming programs go undetected.

\end{frame}
%-----------------------------------------------------------
\begin{frame}\frametitle{Language vs implementation}

The distinction language/implementation is sometimes blurred:

\begin{itemize}
\item Some languages are defined by a \emph{reference
  implementation}.  Examples?
\item Some language definitions are controlled by the same
  organization that supplies some dominating implementation.
  Examples?
\end{itemize}

\end{frame}
%-----------------------------------------------------------
\begin{frame}\frametitle{Language characteristics}

\end{frame}
%-----------------------------------------------------------
\begin{frame}\frametitle{Strong vs weak typing}

A language can be \emph{strongly typed} or \emph{weakly typed}, and
even \emph{untyped}.
\vskip 0.5cm
Strongly typed: It is impossible for an object of one type to be
mistaken for an object of a different type.  Either the compiler made
sure no mistake is possible, or the run-time system checked it.
\vskip 0.5cm
Weakly typed: No such guarantees are made.
\vskip 0.5cm
Untyped: Data can be interpreted differently by different operations.

\end{frame}
%-----------------------------------------------------------
\begin{frame}\frametitle{Static vs dynamic typing}

A language can be \emph{statically typed} or \emph{dynamically
  typed}.
\vskip 0.5cm
Static typing: Type information is associated with the
\emph{variables} in the program.  Type checking or type inference is
always handled at compile time.
\vskip 0.5cm
Dynamic typing: Type information is associated with the \emph{objects}
manipulated by the program.  Type checking must sometimes be done at
execution time.  Not always?

\end{frame}
%-----------------------------------------------------------
\begin{frame}\frametitle{Manifest vs implicit typing}

A statically typed programming language may be based on \emph{manifest
  typing} or \emph{implicit typing}.
\vskip 0.5cm
Manifest typing: The programmers supplies the type of the variables
explicitly.
\vskip 0.5cm
Implicit typing (sometimes called \emph{latent} typing): The compiler
\emph{infers} the type of the variables from the operations it
participates in.

\end{frame}
%-----------------------------------------------------------
\begin{frame}\frametitle{Manifest vs implicit typing}

Why is the distinction between manifest and implicit typing important?
\vskip 0.5cm
With manifest typing, the programmer is given too much responsibility
too early in the development process.  Choices made may easily change
later on.

\end{frame}
%-----------------------------------------------------------
\begin{frame}\frametitle{Manifest vs implicit typing}

Some dynamically typed programming languages allow optional type
declarations.
\vskip 0.5cm
Such declarations are sometimes used in order to allow the compiler to
generate faster code.

\end{frame}
%-----------------------------------------------------------
\begin{frame}\frametitle{Static vs dynamic}

A programming language can be either \emph{static} or \emph{dynamic}.
\vskip 0.5cm
Static: There is a clear distinction between \emph{compile time} and
\emph{run time}.  Code is \emph{generated} at compile time and
\emph{executed} at run time.
\vskip 0.5cm
Dynamic: There is no clear distinction between \emph{compile time} and
\emph{run time}.  The program might change as a result of executing
code (for example, executing a statement that defines a function or a
class).

\end{frame}
%-----------------------------------------------------------
\begin{frame}\frametitle{Interpreted vs compiled}

Can you guess what an \emph{interpreted programming language} is and
what a \emph{compiled programming language} might be?

\end{frame}
%-----------------------------------------------------------
\begin{frame}\frametitle{Interpreted vs compiled}

Can you guess what an \emph{interpreted programming language} is and
what a \emph{compiled programming language} might be?
\vskip 0.5cm
A programming \emph{language} is neither interpreted nor compiled.  A
programming language \emph{implementation} is either one or both (it
is a spectrum of possibilities).

\end{frame}
%-----------------------------------------------------------
\begin{frame}\frametitle{Interpreted vs compiled}

The distinction is important because it affects the \emph{execution
  performance} of the program.
\vskip 0.5cm
An implementation that compiles to native code has the potential of
generating \emph{fast} code.
\vskip 0.5cm
An implementation that has more elements of interpretation can
generate code that is slower by a factor 10 or more compared to code
generated by a native compiler.

\end{frame}
%-----------------------------------------------------------
\begin{frame}\frametitle{Manual vs automatic memory management}

Manual memory management: The language requires the programmer to
de-allocate objects that are no longer going to be referenced.
\vskip 0.5cm
Automatic memory management: The implementation of the language is
required to automatically recycle objects that are no longer
referenced.

\end{frame}
%-----------------------------------------------------------
\begin{frame}\frametitle{Standardization}

A language is said to have an \emph{independent standard} if and only
if the definition of the language is published by an organization
other than a supplier of an implementation.

\end{frame}
%-----------------------------------------------------------
\begin{frame}\frametitle{Scripting languages}

There are few properties that characterize scripting languages.
Probably only:

\begin{itemize}
\item The creators meant for the language to be used for scripting.
\item It is a dynamic language.
\end{itemize}

Common scripting languages are single-implementation languages with
the implementation written as a slow interpreter.  This technique is
considered acceptable because of the first item above.

\end{frame}
%-----------------------------------------------------------
\begin{frame}\frametitle{Scripting languages}

Typically, a static language is used for the main body of code, and a
``scripting language'' for, um, scripting.
\vskip 0.5cm
When the advanced user starts writing serious code using the scripting
language (because that's the only choice possible), the combined
result is slow despite the best intentions of the creators.
\vskip 0.5cm
Furthermore, debugging code written in two languages is typically hard.

\end{frame}
%-----------------------------------------------------------
\begin{frame}\frametitle{Domain-specific languages}

A domain-specific language is a language that was designed for a
particular family of programming tasks.
\vskip 0.5cm
Often, the language is defined by the same organization that then uses
it.
\vskip 0.5cm
The productivity advantage can be huge.
\vskip 0.5cm
Designing and implementing a domain-specific language requires
expertise in language design and compiler technology.

\end{frame}
%-----------------------------------------------------------

\begin{frame}
\frametitle{Thank you}

Questions?
\end{frame}

\end{document}
